\section{Improved closing: $\Theta<\FourRoutesTheta$ with $\vartheta=\FourRoutesDensity$}
\label{sec:four-routes-closing}

This section verifies the price hypotheses of
\cref{thm:density:induction} at
$a=\FourRoutesA$, $K=\FourRoutesK$,
$\vartheta=\FourRoutesDensity$, and thereby proves
\cref{thm:main}. We first account for the normalizing factors of
the positive reference laws, then assemble the improved local bounds,
and finally evaluate the joint comparison and mean inequality.
The quantities $L_0,S,A,D,\mathcal E,T,R_I,J,W$ have the meanings
given in \cref{def:closing-quantities}; their assigned numerical
bounds are collected in a table below.

Two points govern the assembly. First, a middle reference law is
divided by its total mass $Z_i(r)>1$. Its comparison prices carry
this reciprocal, whereas the two positive matching additions retain
their full prices. \Cref{lem:four-routes-normalized} accounts for
both through a common assignment and the reserve $R_Q$. Second,
the higher mean contributions are used inside the joint bound of
\cref{cor:joint-closing}, through the corrections proved in
\cref{subsec:joint-mean}. Their subtraction is justified before
replacing energies by numerical upper bounds.
The density induction uses the wrapping factor
$(R_I/2)K^{-1+\vartheta/36}$.
Appendices~D and~E record the baseline data; the assignments selected
for the final theorem are specified here.

\subsection{Accounting for the normalizing factor $Z_i^{-1}$}

On each middle piece the reference is defined as a sum of records
divided by its total mass $Z_i(r)>1$. The baseline prices of
\cref{def:closing-quantities} were partly stated for the
unnormalized sums; wherever the comparison uses the normalized
laws, the correct price is the baseline price multiplied by
$Z_i(r)^{-1}<1$. This subsection defines an explicit rational
upper bound $c_i$ for $Z_i(r)^{-1}$ on the whole piece, using only
integer floor and ceiling operations, and the following lemma
states the consistent reassignment of all affected prices.
For a prime $b$ with $X_i\le b<Y_i$, put $r=b^{-1/2}$. The mass of the
unnormalized positive reference is
\begin{align*}
 Z_i(r)={}&1+(3w_{{\rm D},i}+2w_{{\rm S},i}+v_{{\rm L},i})r
 +(6w_{{\rm R},i}+3v_{{\rm D},i}+2v_{{\rm S},i})r^2\\
 &+(2w_{{\rm C},i}+6v_{{\rm R},i})r^3+2v_{{\rm C},i}r^4.
\end{align*}
All coefficients are nonnegative. The following explicit rational numbers
therefore provide simultaneous majorants on the entire piece:
\begin{equation}\label{eq:four-routes-normalizer}
 r_i^-={\left\lfloor\sqrt{\lfloor10^{120}/Y_i\rfloor}\right\rfloor\over10^{60}},
 \qquad c_i=10^{-30}\left\lceil {10^{30}\over Z_i(r_i^-)}\right\rceil,
 \qquad Z_i(r)^{-1}\le c_i<1.
\end{equation}
The floor and ceiling are integer operations. In particular no numerical
square-root approximation is needed to define $c_i$.

\begin{lemma}[Consistent normalized prices]\label{lem:four-routes-normalized}
Define new assigned prices from those of \cref{def:closing-quantities} as follows.
Every norm price and raw-arm price in the middle row
\eqref{eq:middle-piece-prices} is multiplied by $c_i$. For each first-stage
total $T\in\{U_0',P',D',V'\}$, let $t_i$ be its exact certified upper
charge in piece $i$, and replace $T$ by
\begin{equation}\label{eq:four-routes-first}
 T^{\rm new}=\left(T-\sum_i t_i\right)+\sum_i c_i t_i.
\end{equation}
The first parenthesis is its retained nonnegative rounding slack.
Replace the invariant-background assignment by
${\Uinv}^{\rm new}=\sum_i c_i\Uinv_i$ and the one-tag symmetric
assignment by $s^{\rm new}=\sum_i c_i s_i$, where $s_i$ is the assigned
symmetric folded-tag price in that piece. Use
\[
 \delta_{\rm s}^{\rm new}=D^{\prime\rm new}/2-s^{\rm new}
\]
in every orbit saving, and recompute the full tagged series and every
class polynomial from these assignments.

Keep the chosen $F$ matching-side price $T_K$ unchanged. If
$S_{E,i}^{\rm old}$ and $A_{E,i}^{\rm old}$ are the unnormalized
assigned matching prices, use
\begin{equation}\label{eq:four-routes-q-reserve}
 T_{QE}=t\sum_i(S_{E,i}^{\rm old}+A_{E,i}^{\rm old})
       =S_E^{\rm new}+A_E^{\rm new}+R_Q,
 \quad
 R_Q=t\sum_i(1-c_i)(S_{E,i}^{\rm old}+A_{E,i}^{\rm old}),
\end{equation}
where $a$ is the target exponent, $a_0$ is the middle anchor exponent, and $t=e^{2(a-a_0)570}$. Thus $R_Q$ is added to the closing quantity $T$ of \cref{def:closing-quantities} and to the exponent in the mean-error factor $R_{\rm err}e^{T_K+T_{QE}}$. These assignments preserve all
comparison, class and matching inequalities.
\end{lemma}

\begin{proof}
Each middle comparison difference, each of its arm coefficients, and
each first-stage record has the common physical factor $Z_i(r)^{-1}$.
Its former assigned price omitted this factor. Equation
\eqref{eq:four-routes-normalizer} consequently bounds every term after
multiplication by $c_i$. This is a simultaneous coefficientwise bound
on the full positive majorant polynomial; the leading, remainder and
full arm prices are changed together. The assigned slack in
\eqref{eq:four-routes-first} is nonnegative and remains available.
The two folded tags have equal coefficients, and retain equal sharp
and variation prices after this operation.

The invariant and generic pieces of the first-stage bracket are changed
with the same factors. The exact checks below give
$0\le{\Uinv}^{\rm new}\le U_0^{\prime\rm new}$ and
$0\le s^{\rm new}\le D^{\prime\rm new}/2$. Consequently the
invariant-background and one-delta reductions are the same algebraic
reductions of this new positive majorant as in
\cref{def:closing-quantities}(6),(7). Keeping their former, larger
coupons while lowering only the generic total would not justify the
subtraction; the simultaneous assignments avoid that error.

The physical reference and its actual half marginals are unchanged, as are the quantities entering restoration, namely the likelihood and chi-square bounds, the piece product bounds and the chosen matching-side coefficients. In particular, the chosen $K_{\rm L},K_{\rm R}$
and $Q_E$ of \cref{sec:matching-mean} are unnormalized positive
additions. Their former prices remain valid. Equation
\eqref{eq:four-routes-q-reserve} charges the entire old $Q_E$ price
while using the smaller comparison price elsewhere. The old direct
$F$ matching price and its positive increments are also retained.
This proves the matching and mean-error assertions. All restoration
and class gains are then evaluated using the same new price assignment,
so no term is counted using incompatible generic and discounted prices.
\end{proof}

\begin{lemma}[Upper-band normalization]\label{lem:four-routes-upper-normalized}
For a band $[X_i,Y_i)$, let $P_{\mu,i}$ be the positive mass polynomial
of \cref{def:upper-reference}. Write
$d_{\rm L},d_{\rm D},d_{\rm S},d_{\rm R},d_{\rm C}=0,1,1,2,3$ and
$m_{\rm L},m_{\rm D},m_{\rm S},m_{\rm R},m_{\rm C}=1,3,2,6,2$.
Thus, for the $25$ ordered states $uv$ with $u,v\in\{\mathrm L,\mathrm D,\mathrm S,\mathrm R,\mathrm C\}$,
\[
 P_{\mu,i}(r)=\sum_{u,v}m_um_v
       \bigl(w_{uv}r^{d_u+d_v}+v_{uv}r^{d_u+d_v+1}\bigr).
\]
Define
\[
 r_i^-=\frac{\left\lfloor\sqrt{\lfloor10^{200}/Y_i\rfloor}\right\rfloor}{10^{100}},
 \qquad c_i^{\rm u}=10^{-30}\left\lceil\frac{10^{30}}{P_{\mu,i}(r_i^-)}\right\rceil.
\]
Then $P_{\mu,i}(b^{-1/2})^{-1}\le c_i^{\rm u}\le1$.
Multiplying each band's assigned symmetric and skew prices by
$c_i^{\rm u}$ also gives valid, simultaneously smaller generic matching
prices. The tail prices and chosen middle matching diagonals are unchanged.
\end{lemma}

\begin{proof}
The reference polynomial has constant term one and nonnegative
coefficients. The same monotonicity argument as in
\eqref{eq:four-routes-normalizer} proves the stated bound. In the proof
of \cref{prop:record:upper}, every repair, deletion and skew envelope
is divided by this common polynomial before that denominator is
discarded in the displayed prices. Retaining its reciprocal upper
bound is therefore valid for every positive coefficient of the three
comparison polynomials. The upper generic matching envelope is dominated by the sum of their symmetric and skew masses. The same finite upper diagonal assignments are used in the joint mean
allocation of \cref{prop:joint-upper-mean}. This normalization uses the
weights already selected in \cref{subsec:upper-cross-weighted};
it does not make any further change of reference law. The common
factor cancels in the packet mixture gates.
\end{proof}

\subsection{The selected laws and common assignments}

Fix
\[
 a=\FourRoutesA,\qquad \vartheta=\FourRoutesDensity,\qquad
 K=\FourRoutesK,\qquad
 \eps=\frac{a}{72-\vartheta}=\FourRoutesExponent.
\]
The $72$ slots, six queries, $33$ middle pieces, cutoff $Q_3=19951586$,
upper endpoint $P_3=167\cdot2^{40}$, $310$ finite upper bands,
local coefficient bound $585$, block budget $570$, and external depth
$64$ remain as specified in the construction. The complete middle
portfolio is the one specified in
\cref{sec:middle-reference,sec:positive-roots}.

Use the independent-draw prefix of \cref{thm:density:iid-prefix}, the
accepting middle constants and merges of
\cref{lem:merge-repair,lem:record:accepting-c2}, and the new finite upper
comparisons of \cref{subsec:upper-six-refinements}, with the complete
upper reference reweighting of \cref{subsec:upper-cross-weighted}.
Each component keeps its original target marginals. The density-moment
bounds apply componentwise and hence to these positive mixtures.
The finite upper changes preserve acceptance of all six queries.

The middle charge anchor is $a_0=\FourRoutesAnchor$. Its exact
histograms and analytic prime-counting cells give charges of degrees
$2,3,4$. Reprice every middle charge by
$t=e^{2(a-a_0)570}$; this uses the certified maximum block product,
not an average length. Each upper term retains its own certified
product-log bound and anchor. Multiply the middle and upper
assignments by their normalizing factors, using the new full upper
reference polynomial, as in
\cref{lem:four-routes-normalized,lem:four-routes-upper-normalized}.
For each affected upper state replace its entire symmetric and skew
coefficient by the corresponding aggregate cap in
\cref{prop:upper-soft-nine,prop:upper-ss-six,prop:upper-constant-six};
for LR and RL subtract exactly the rebates in
\cref{prop:upper-observed-dense}. Every affected charge is then
resummed. The other upper coefficients retain
\cref{lem:upper:neumann-refinement}. This changes $S,A,T$ and all
uses of those totals together.

Every restoration is recomputed at $a$ with the density-valid bound
\begin{equation}\label{eq:four-routes-restoration}
 \mathcal F_2(C,h,M,\chi)=
 \max\left\{0,\frac{\ell}{M},
 \frac{\ell-\sqrt{\chi((1+\chi)C^2-\ell^2)}}{1+\chi}\right\},
 \qquad \ell=\max\{0,1-(C-1)h\},
\end{equation}
proved in \cref{sec:density-upgrade}. This includes the selected,
background and pooled variance gains. All sixteen middle envelope
classes, their invariant refinements, and their orbit coupons use
the same normalized raw-arm and outside assignments.

For the mean keep only the terms supported on one block in
\cref{prop:record:mean-credit}, whose sum is denoted $m_{\rm nat,1}$.
Replace that proposition's higher-degree mean credits by
\cref{prop:joint-middle-mean,prop:joint-upper-mean}, whose four
middle families and nonempty upper family are disjoint. Set
\[
 m_1=1+m_{\rm nat,1}-R_{\rm err}e^{T_K+T_{QE}},\qquad
 \Delta_{\rm mid}=\Delta_\alpha+\Delta_M+\Delta_O+\Delta_K.
\]
The separate chosen-$Q_E$ reserve of
\eqref{eq:four-routes-q-reserve} is included in both $T$ and the mean
error exponent. The upper-triggered allocation uses only the finite
upper assignments; all prefix and tail padding remains in the positive
costs, matching and wrap.

\subsection{The joint closing number}

The finite conductor moment is evaluated at the selected $(a,K)$
using \cref{prop:imp:padic-price,cor:imp:padic-majorant}. Put
\begin{align*}
 \mathcal E&=B(e^S\cosh A-1)-D,&
 J&=Be^T(R_I/2)^6,\\
 W&=Be^{S+A}\bigl[(1+(R_I/2)K^{-1+\vartheta/36})^6-1\bigr].
\end{align*}
The joint comparison and mean bound proved above gives
\begin{equation}\label{eq:final-joint-closing}
 \Theta=\frac{63}{64}(L_0+\mathcal E)+J+W+
 \frac{1-m_1}{64}-\Delta_{\rm mid}-\Delta_u+\mathcal R_u.
\end{equation}
The two negative allocations are justified inside the common
positive completion. They are not additional independent instances
of (CP7). This is the joint criterion of \cref{cor:joint-closing}.

\begin{table}[ht]
\centering
\caption{The improved closing ledger. Bounds are rounded outward from the
complete enclosures; the closing uses the enclosures themselves.}
\label{tab:four-routes-ledger}
\begin{tabular}{@{}lr@{\qquad}lr@{}}
\toprule
quantity & bound & quantity & bound\\
\midrule
$B$ & $<2.373803797078$ & $L_0$ & $<0.000507024733$\\
$S$ & $<0.008981913739$ & $A$ & $<0.985245889185$\\
$T$ & $<0.610594125990$ & $D$ & $>0.495667577473$\\
$\mathcal E$ & $<0.785417906625$ & $m_1$ & $>1.529526120248$\\
$R_I$ & $<1.252937308029$ & $J$ & $<0.264252164235$\\
$W$ & $<0.000041043577$ & $D_{\rm var}$ & $>0.009861452145$\\
$\Delta_\alpha$ & $>0.007949257298$ & $\Delta_M$ & $>0.013289595166$\\
$\Delta_O$ & $>0.006920381923$ & $\Delta_K$ & $>0.000788314597$\\
$\Delta_u$ & $>0.000716967431$ & $\mathcal R_u$ & $<0.000000070807$\\
\bottomrule
\end{tabular}
\end{table}


\begin{certificate}[The final joint closing]\label{cert:four-routes-closing}
The construction just specified has
\[
 \Theta<\FourRoutesTheta<1,\qquad
 1-\Theta>2.29\cdot10^{-7}.
\]
The reproducible entry point is
\path{research/epsilon_2026_10_29/verify.py}.
Its author assembly and independent complete assembly give intersecting
outward enclosures with value
$0.9999997708753344\ldots$.
The independent assembly reconstructs all $33$ middle price rows,
$310$ upper normalizers and comparison polynomials, all restoration
calls, the chosen-matching reserve, the finite conductor moment,
the physical resource ratios, and the $192$ shared residual gates.
It computes \eqref{eq:final-joint-closing} directly. Both assemblies
use the same completed reference reweighting; its independent local
review checks $602$ root bounds and $32\,550$ completion inequalities.

The new local checks independently reconstruct all sixteen soft
comparison graphs and complete character partitions, their exact
leading matrices, scalar masses, count orders, repairs and tensor
gates; all fourteen dense law/mixture matrix tests and emission
rebates; and the complete SS and LC/CL scalar proofs. The exact
analytic, prefix and density gates hold on $0<a\le91/100000$.
The local and global records, with source hashes, are in the
\path{certificates/} and \path{verification/} subdirectories of
\path{research/epsilon_2026_10_29/}. Numerical verification establishes these finite
assertions; the positive-energy allocation and operator reductions
are the proofs in the preceding sections.
\end{certificate}

\begin{proof}[Completion of the proof of \cref{thm:main}]
The actual reference and its retained marginals satisfy the moment
catalogue and packet bounds of \cref{sec:density-upgrade} at
$\vartheta=\FourRoutesDensity$. The new accepting upper laws are
repaired to those same marginals, so their positive diagonal
reductions use the same catalogue. All generic selected packets
have the required unit mixed moment; the stronger weighted
restoration moments are checked for the actual middle fold laws.
The first-stage assignments, chosen matching prices, and complete
upper comparisons are compatible by the preceding construction.

The joint mean proof has strictly positive residual coefficients.
Consequently its completed positive series has the same robust
inflation and the same two analytic error types as the original
closing proof. The induction of \cref{thm:density:induction}, in
its joint form, therefore applies. Its extended analytic domain is
valid: the relative-error margin at $a=91/100000$ is
$943271/3284710000>0$, and the prefix plus tail is below the assigned
$10^{-12}$ in every channel. The strict closing is
\cref{cert:four-routes-closing}.

Initialization at the finitely many small scales and continuity in
the inflation parameter give, for every square-difference-free
$A\subseteq[\Kint\Lsc]$ of density $\alpha$,
\[
 Y(\Lsc,A)\le C\Lsc^{-a}\alpha^\vartheta.
\]
Reflection positivity gives $\alpha^{72}\le Y(\Lsc,A)$; for $\alpha>0$,
division and the $(72-\vartheta)$th root yield
$\alpha\le C'\Lsc^{-a/(72-\vartheta)}$.
The fixed relation $N=K\Lsc$ and interval rounding change only the
absolute constant. Hence $s(N)\le C''N^{1-\FourRoutesExponent}$
for every $N\ge1$.
\end{proof}
